% 00-abstract.tex -- placed inside \begin{abstract}...\end{abstract} by main.tex.
% Plain prose only (it is also pasted into HotCRP). Scenario text uses \armswitch.
% Wording rules (brief, part d): keep "introducing ... into the model review
% process" until Section VII documents a real review (D16, N8); drop "an
% industrial case study" if Section VII degrades to a retrospective case; in the
% withdrawn scenario, drop the fairness-test clause if N2 did not run and the
% test is not stated as a design contribution in Section VI-B.
European banks put credit-scoring models through formal review, and the EU AI Act classifies creditworthiness assessment as a high-risk use.
The same control functions now receive large language models, for which model-risk practice has no established review procedure.
We report on introducing VERA into a European bank's model review process.
VERA is an evaluation engine that turns a declared specification of responsible-AI requirements into scored, traceable evidence.
Its central design decision keeps that specification, scoring policy included, as data outside the engine core.
\armswitch{%
We test this decision on two model families, each with its own specification and reviewer population: cloud-served language models, examined with the team that assesses generative-AI use cases, and tabular credit-scoring classifiers, examined with the legal and risk function.
Fairness is the requirement both populations own.
Because common fairness criteria cannot in general hold together, we treat the criterion as an explicit, auditable policy choice and identify which review verdicts survive a change of criterion.
}{%
We test this decision on two model families, each with its own specification and reviewer population: cloud-served language models, examined with the team that assesses generative-AI use cases, and, as a portability demonstration, a tabular credit-scoring classifier examined with the legal and risk function.
Fairness is the requirement both populations own.
Because common fairness criteria cannot in general hold together, we treat the criterion as an explicit, auditable policy choice and identify which review verdicts survive a change of criterion.
}{%
We apply VERA to cloud-served language models examined by two populations, the team that assesses generative-AI use cases and the legal and risk function, and specify how the same engine accepts tabular credit-scoring classifiers.
Because common fairness criteria cannot in general hold together, we treat the criterion as an explicit, auditable policy choice and give a test of which verdicts survive a change of criterion.
}%
We also map, by serving mode and model family, which requirements are measurable natively, through a fallback or not at all, and report a decision study with both populations, an industrial case study and lessons for practice.
